\chapter{Blepharoplasty}

\section*{Overview}
Blepharoplasty is surgery of the eyelids that removes or repositions excess skin, muscle, or fat to improve eyelid function, visual-field obstruction, or appearance.

\section*{Why It Is Done}
It may be performed for upper-eyelid skin that obstructs vision, lower-eyelid bags, eyelid asymmetry, or cosmetic concerns after appropriate eye evaluation.

\section*{How It Is Performed}
Through carefully placed eyelid incisions, the surgeon removes or redistributes tissue and closes the skin. Local anesthesia with sedation or general anesthesia may be used.

\section*{Risks and Recovery}
Bruising, swelling, dry eyes, asymmetry, infection, bleeding, scarring, difficulty closing the eyes, and rare vision-threatening complications are possible. Cold compresses, head elevation, and follow-up protect healing.

\section*{Outcome and Follow-up}
The expected outcome depends on the reason for the procedure, the person's health, and whether additional treatment is needed. The clinician reviews the result with the patient, explains any next steps, and sets follow-up according to the findings and recovery.

\section*{When to Seek Medical Care}
Follow the procedure team's specific instructions. Seek urgent medical assessment for severe or worsening pain, uncontrolled bleeding, fever or signs of infection, trouble breathing, chest pain, fainting, new weakness or confusion, inability to urinate, or any rapidly worsening symptom. The appropriate warning signs depend on the procedure and the patient's medical condition.

\section*{References}
\begin{enumerate}
  \item MedlinePlus. Eyelid Surgery. U.S. National Library of Medicine; 2025.
  \item Current Medical Diagnosis and Treatment. McGraw-Hill; 2025.
\end{enumerate}
\clearpage
